\section{Conditional laws and the Poisson target}\label{supp:sec:conditional-details}
The arithmetic theorems provide total-variation errors. This appendix
records exactly when those errors are small enough to transfer a local
probability, a conditional state, or a Gaussian limit. Properties of the
independent Poisson target are not additional arithmetic results.

\subsection{Scalar Gaussian and local estimates}
Let $P_\lambda\sim\Pois(\lambda)$ and let $\epsilon_N$ denote any proved
conditional scalar total-variation bound in the article. If $\lambda\to\infty$,
\[
 \sup_t\left|\PP\left(\frac{Z-\lambda}{\sqrt\lambda}\le t\mid C\right)
                       -\Phi(t)\right|
 \le\epsilon_N+C\lambda^{-1/2}.
\]
For example, write $P_\lambda$ as a sum of $\lceil\lambda\rceil$
independent Poisson variables of mean at most one and use the classical
Berry bound \cite{Berry1941}. Their third absolute centered moments
are $O(\lambda/\lceil\lambda\rceil)$, which gives the displayed scale.
The normalization uses the deterministic comparison mean $\lambda$;
replacing it by an actual conditional mean requires a separate moment
estimate and is not inferred from total variation alone.

For every integer $n\ge1$ and every $\lambda>0$, one has the exact form
\begin{equation}\label{supp:eq:poisson-atom}
 \PP(P_\lambda=n)=\frac{e^{-\delta_n}}{\sqrt{2\pi n}}
                    \exp[-\lambda h(n/\lambda)],\qquad
 0\le\delta_n\le\frac1{12n}.
\end{equation}
Here $h(u)=u\log u-u+1$, and $\delta_n$ is independent of $\lambda$. To check the remainder, set
\[
 \delta_n=\log(n!)-(n+\tfrac12)\log n+n-\tfrac12\log(2\pi).
\]
Stirling's limit gives $\delta_n\to0$. With $t=(2n+1)^{-1}$,
\[
 \begin{aligned}
 \delta_n-\delta_{n+1}
 &=(n+\tfrac12)\log(1+1/n)-1
   =\sum_{j\ge1}\frac{t^{2j}}{2j+1},\\
 0\le\delta_n-\delta_{n+1}
 &\le\frac{t^2}{3(1-t^2)}=\frac1{12n(n+1)}.
 \end{aligned}
\]
Telescoping to infinity proves the stated bounds. Substitution of the
factorial identity into $e^{-\lambda}\lambda^n/n!$ proves
\eqref{supp:eq:poisson-atom}; moreover
$0\le1-e^{-\delta_n}\le1/(12n)$, uniformly in the intensity.
For $n=0$ the mass is exactly $e^{-\lambda}$ and the square-root formula
is not used.

Thus relative transfer at $n$ requires
$\epsilon_N/\PP(P_\lambda=n)\to0$. For the hard scalar rate a sufficient
condition is
\[
 I+\log^+\lambda+\tfrac12\log n+\lambda h(n/\lambda)
                  \le V_N-c\nu_N,
\]
and for the soft rate a sufficient condition is
\[
 2I+\log^+\lambda+\log n+2\lambda h(n/\lambda)
                  \le\widehat V_N-c\widehat\nu_N.
\]
In both cases the event must be measurable at the corresponding prime
cutoff. The logarithms of the reciprocal target masses are then
$O(V_N)$ or $O(\widehat V_N)$, so division by that mass does not destroy
the polynomial remainder.

In the central range
$n=\lambda+t\sqrt\lambda+O(1)$, $|t|=o(\lambda^{1/6})$,
Taylor expansion gives
$\lambda h(n/\lambda)=t^2/2+o(1)$ and $\log n=\log\lambda+o(1)$.
The respective central local budgets are therefore
\[
 I+\tfrac32\log\lambda+\tfrac12t^2\le V_N-c\nu_N,
 \qquad
 2I+2\log\lambda+t^2\le\widehat V_N-c\widehat\nu_N.
\]
A global scalar approximation at its upper intensity endpoint does not
by itself imply a relative local approximation there.

For completeness, the analogous relative upper-tail transfer has the
additional denominator $\PP(P_\lambda\ge\lceil\lambda+t\sqrt\lambda\rceil)$.
Stirling's formula and summation of successive Poisson masses yield the
usual Gaussian moderate-tail comparison for
$0\le t=o(\lambda^{1/6})$. Equivalently, expanding the exponent uniformly
on the contributing $O(\sqrt\lambda/(1+t))$ scale gives relative error
$O((1+t^3)/\sqrt\lambda)$; the exponentially smaller outer part of the
sum is bounded by the ratio of successive masses. Since
$1-\Phi(t)\asymp e^{-t^2/2}/(1+t)$ for $t\ge0$, a sufficient hard
arithmetic condition is
\[
 I+\log^+\lambda+\tfrac12t^2+\log(1+t)\le V_N-c\nu_N.
\]
This condition must be combined with the intrinsic Poisson restriction
$t=o(\lambda^{1/6})$. It is not a claim about larger deviation scales.

\subsection{Resolving an entire conditional path}
Let $\epsilon^{\rm path}$ be the aggregate staircase error based at the
present threshold. The sharpened conditioning inequality
\eqref{eq:sharp-conditioning} gives a full future-path error at most
$\epsilon^{\rm path}/p_\lambda(n)$, where
$p_\lambda(n)=\PP(P_\lambda=n)>\epsilon^{\rm path}$.
The same assertion holds jointly with any fixed available lower segment.
Forward transitions are binomial thinning; backward transitions add
independent Poisson increments.

Before division, the aggregate proof gives
\[
 \epsilon^{\rm path}\ll
 e^I\lambda(1+\log^+(2\lambda))e^{-V_N+o(\nu_N)}
                      +N^{-1/3+\varepsilon},
\]
provided the information, log-intensity and mark cutoff are $O(V_N)$.
Combining this with \eqref{supp:eq:poisson-atom} proves the budget
\eqref{eq:resolved-budget}. Its hypotheses also validate that use of the
common logarithmic band. Indeed, for $n\ge2$, if $\lambda\ge1$ the budget
bounds $\log\lambda$ by $O(V_N)$. If $\lambda<1$, use
\[
 \lambda h(n/\lambda)=n\log n-n+\lambda-n\log\lambda
\]
to conclude $|\log\lambda|=O(V_N)$. Thus the base offset, information
cost and chosen mark cutoff are all $O(V_N)=o(\log N)$.
The atom's reciprocal is $\exp(O(V_N))=N^{o(1)}$, which can be absorbed
by first choosing a smaller arithmetic $\varepsilon$.
For a fixed present state at a fixed positive limiting intensity the
atom stays bounded away from zero, and no extra logarithmic cost is needed.

The condition $p_\lambda(n)>\epsilon^{\rm path}$ only makes the
conditional law well-defined and gives a bound. Convergence requires
$\epsilon^{\rm path}/p_\lambda(n)\to0$. Nothing here supplies a single
finite-intensity comparison for the entire reverse path: its target
mass diverges toward thresholds $m\to-\infty$.

\subsection{Quenched transfer along geometric scales}
The proof of \cref{cor:quenched} gives the single-length
Markov--Borel--Cantelli argument. Its uniform extension uses the bound
\[
 e^{-(c-\beta)\nu_{N_k}+o(\nu_{N_k})}
       +e^{\beta\nu_{N_k}}N_k^{-1/3+\varepsilon},\qquad 0<\beta<c,
\]
for the exceptional probability on $N_k\asymp q^k$, $q>1$.
Since $\nu_{N_k}\asymp\sqrt{k/\log k}$, multiplication by
$O(\log N_k)=O(k)$ preserves summability. A single union bound therefore
covers every integer length in a fixed logarithmic band, provided the
error margin is uniform. The resulting almost-sure statement concerns
these conditional laws at individual scales; it requires no independence
between scales and asserts no canonical coupling over all scales.

\subsection{Exact target consequences}\label{supp:sec:target-consequences}
For a phase $\vartheta\in[0,1]$, let
$P_r\sim\Pois(2^{\vartheta-r-1})$ independently for $r\in\Z$, and
$S_m=\sum_{r\ge m}P_r$. Each tail sum is finite almost surely. Then
\[
 S_m\sim\Pois(2^{\vartheta-m}),\qquad
 \operatorname{Cov}(S_m,S_n)=2^{\vartheta-\max(m,n)}.
\]
Splitting $S_m=P_m+S_{m+1}$ gives the exact forward thinning and
backward immigration kernels in the article. For $s\ge0$ the joint
probability generating function is
\[
 \EE\prod_{j=0}^s z_j^{S_{m+j}}
 =\exp\left[
 \sum_{r=m}^{m+s-1}2^{\vartheta-r-1}
       \left(\prod_{j=0}^{r-m}z_j-1\right)
 +2^{\vartheta-m-s}\left(\prod_{j=0}^sz_j-1\right)\right].
\]
Conditionally on $S_m=n$, the $n$ points have independent uniform
positions and independent geometric overshoots
$\PP(E=e)=2^{-e-1}$. For lattice targets the positions are uniform on
the corresponding finite grid. Conditioning the critical prelimit field
on a fixed count transfers this assertion by the sharp conditioning
lemma; only the final grid-to-uniform passage is weak convergence.

The complete centered target is locally finite on $[1,2]\times\Z$,
although it has infinite total mass toward $r=-\infty$. For every
nonnegative continuous compactly supported $g$ put
\[
 \Psi_\vartheta(g)=
 \exp\left[-\int_1^2\sum_{r\in\Z}2^{\vartheta-r-1}
                                (1-e^{-g(t,r)})\,dt\right].
\]
Compact support meets finitely many levels. At those levels the lattice
comparison and a Riemann sum give
\[
 \EE(e^{-\langle g,\mathcal E_N\rangle}\mid C_N)
                         =\Psi_{\vartheta_N}(g)+o_g(1)
\]
whenever $I_N\le V_N-c\nu_N$. A convergent phase subsequence is needed
only to replace this moving deterministic target by a single limit law.

For a base intensity $\lambda_N\to\infty$, the exact maximum identity
in the article, applied at
$m_N=\lfloor\log_2\lambda_N\rfloor+j$, gives, along
$\{\log_2\lambda_N\}\to\vartheta$,
\[
 \PP(M_{N,L}-\lfloor\log_2\lambda_N\rfloor\le j)
                         \longrightarrow e^{-2^{\vartheta-j-1}}.
\]
This is a lattice extreme law. The identity uses the scalar comparison
at the shifted critical length, not an unproved uniform law for every
large marked field.

Finally, the Poisson--Gaussian bridge follows by splitting the independent
levels at a fixed critical level $r_0$. Below it, the normalized Poisson
sums satisfy a multivariate central limit theorem with covariance
$2^{-\max(j,k)}$. The centered contribution of the levels at or above
$r_0$ has variance $O(\Lambda_N^{-1})$ after normalization, whereas
those levels retain their Poisson law and are independent of the lower
part. This proves joint independence in the limit. Continuous-time
embeddings and phase-removing auxiliary randomizations belong to the
separate process article \cite{PoulyLatticePoissonFlows2026}; no such enlargement
is used in the arithmetic proofs here.
